% Section 8.4, from lectures/L08/appendix/c_exercises.md.

\section{Exercises}
\label{c8:sec:exercises}\label{c8:app:c}

\begin{aside}
Nine, ending with the Cross-check. Do them in order; Exercise~\ref{c8:ex:6} onwards need the target
running.

Where an exercise can be checked mechanically, a \textbf{Check yourself} line says how.

\textbf{A warning about Exercise~\ref{c8:ex:8}.} It asks you to load a handler that never
acknowledges its device. That wedges the machine, on purpose, and the only way out is a reboot. A
few seconds here; do not run the equivalent against anything you cannot power-cycle.

\textbf{Solutions.} The exercises that are not Code are answered in \secref{sol:sec:c8}. The Code
exercises have no solution: each is checked by running it, as its Check yourself line says.
\end{aside}

% ------------------------------------------------------------------------------------------
\exercise[fillaccent2]{The three numbers}{Recall}
\label{c8:ex:1}

\xpart{a} The device tree says \code{interrupts = <0 112 4>} and \file{/proc/interrupts} says
\code{17 ... GIC-0 144}. Explain all three numbers and where each comes from.

\xpart{b} $112 + 32 = 144$. What is the 32?

\xpart{c} Which of the three changes if you boot the same kernel on the same machine with one extra
device present? Why?

\xpart{d} A datasheet documents the device's interrupt as 112. A colleague greps
\file{/proc/interrupts} for 112, finds nothing, and concludes the driver is not loaded. What has
gone wrong with their reasoning?

\xpart{e} What does the third cell, \code{4}, mean, and what would \code{1} have meant instead?

% ------------------------------------------------------------------------------------------
\exercise[fillaccent2]{Reading \code{/proc/interrupts}}{Recall}
\label{c8:ex:2}

\begin{console}
 17:       1789          0  GIC-0 144 Level     qa_irq
\end{console}

\xpart{a} Name every field.

\xpart{b} The second CPU's count is 0. Give two different explanations, and say how you would tell
them apart.

\xpart{c} The count is not moving and the device is not working. What have you ruled out, and what
is still possible?

\xpart{d} The count is increasing by millions per second. What is almost certainly wrong?

\xpart{e} You unload your module and load it again. Does the count restart from zero? What does that
imply for any measurement of an interrupt rate?

% ------------------------------------------------------------------------------------------
\exercise[fillaccent3]{The acknowledgement}{Hand calculation}
\label{c8:ex:3}

\code{qa-dev}'s \code{IRQ_STATUS} is write-one-to-clear. At the moment your handler runs, it holds
\code{0b01} (\code{QA_DEV_IRQ_SAMPLE}). For \textbf{c)} to \textbf{e)}, suppose the bit is cleared
only by a write to it; \textbf{f)} is about the other way this device has.

\xpart{a} Your handler writes \code{0xFFFFFFFF} to it. What happens to the \code{SAMPLE} bit? What
happens to an \code{OVERRUN} that the device set one microsecond after your read?

\xpart{b} Your handler writes back exactly the value it read. Same two questions.

\xpart{c} Your handler acknowledges first and then drains the FIFO. Describe the sequence of events
that produces a spurious interrupt, and say whether any data is lost.

\xpart{d} Your handler drains the FIFO and then acknowledges with a value it read \emph{before}
draining. Describe the sequence that loses an event.

\xpart{e} Of \textbf{c)} and \textbf{d)}, one is survivable and one is not. Which, and why?

\xpart{f} For this device there is a second way to clear \code{IRQ_SAMPLE} that involves no write to
\code{IRQ_STATUS} at all. What is it, and how would you find that out from
\file{tools/qa-dev/qa-dev.c}?

% ------------------------------------------------------------------------------------------
\exercise[fillaccent3]{An interrupt budget}{Hand calculation}
\label{c8:ex:4}

A device raises an interrupt every 250 microseconds. Its handler takes 8 microseconds.

\xpart{a} How many interrupts per second, and what fraction of one CPU does the handler consume?

\xpart{b} The handler grows to 40 microseconds. Recompute both.

\xpart{c} At what handler duration does the CPU do nothing but service this device?

\xpart{d} The device's FIFO holds 16 samples and it produces one per interrupt. How long can the
system ignore the device before samples are dropped?

\xpart{e} You move 30 of the 40 microseconds into a threaded handler. Which of your answers change,
and which does not? In particular, does the \emph{system's} total work change?

% ------------------------------------------------------------------------------------------
\exercise[fillaccent4]{Where does the work go}{Design}
\label{c8:ex:5}

For each, choose hard handler, threaded handler, or workqueue, and justify it in one sentence.

\xpart{a} Reading a status register and clearing an interrupt.

\xpart{b} Copying 4 KB out of the device's FIFO into a driver buffer.

\xpart{c} Waking a userspace reader that is blocked on \code{read()}.

\xpart{d} Allocating a buffer with \code{GFP_KERNEL}.

\xpart{e} Retrying a failed transfer after a 10 ms delay.

\xpart{f} Updating a statistics counter.

\xpart{g} Logging an error with \code{dev_err}.

Then: \textbf{h)} two of your answers would be illegal in a hard handler. Which, and what exactly
goes wrong? \textbf{i)} Which config option this course enables would tell you?

% ------------------------------------------------------------------------------------------
\exercise{Take the interrupt}{Code}
\label{c8:ex:6}

Write \file{lectures/L08/lab/qa_irq.c}. It must:
\begin{itemize}
  \item Find the device's node with \code{of_find_compatible_node} using
    \code{QA_DEV_DT_COMPATIBLE}, get its address with \code{of_address_to_resource}, and its
    interrupt with \code{irq_of_parse_and_map}. Print both.
  \item Claim and map the registers as in Chapter~\ref{c6:ch}, and check the identity register.
  \item Register a handler with \code{request_irq}, named \code{qa_irq}.
  \item Program \code{PERIOD_NS} from a module parameter (default 1000000) and set \code{ENABLE |
    IRQ_EN}.
  \item In the handler: read \code{IRQ_STATUS}, \textbf{return \code{IRQ_NONE} if nothing is
    pending}, acknowledge with the value you read, then drain the FIFO counting samples.
  \item Count hard interrupts and samples in \code{atomic_t}s and print them on unload.
  \item Take a \code{threaded} module parameter which, when set, uses \code{request_threaded_irq}
    with \code{IRQF_ONESHOT} and does the draining in the threaded half.
  \item Disable the device and \code{free_irq} on unload, in that order.
\end{itemize}

\checkyourself \code{make test L=L08} reports \textbf{PASSED}, with nine checks.
\file{/proc/interrupts} gains a \code{qa_irq} line whose count climbs at roughly a thousand per
second.

\xpart{a} Why must the device be disabled before \code{free_irq}, and not after?

\xpart{b} Your handler returns \code{IRQ_HANDLED} unconditionally instead of checking. On this
target nothing appears to break. Explain what you have broken anyway, and on what kind of system it
would show up.

% ------------------------------------------------------------------------------------------
\exercise{The threaded half}{Code}
\label{c8:ex:7}

With \code{threaded=1}:

\xpart{a} Find the kernel thread with \code{ps}. What is it called, and what does the name tell you?

\xpart{b} The hard handler now returns \code{IRQ_WAKE_THREAD}. Where must the acknowledgement
happen, and what happens if you leave it in the threaded half without \code{IRQF_ONESHOT}?

\xpart{c} Measure the interrupt rate both ways over two seconds. Is the threaded version slower? Was
that your prediction?

\xpart{d} Name two things the threaded handler can do that the hard one cannot, and one thing it
gives you that a workqueue does not.

% ------------------------------------------------------------------------------------------
\exercise[fillaccent4]{The interrupt that never stops}{Design}
\label{c8:ex:8}

\textbf{Read the warning at the top of this section.}

\xpart{a} Remove everything in your handler that clears \code{IRQ_SAMPLE}: the acknowledgement, and
the drain with it, because emptying the FIFO clears the bit too (Exercise~\ref{c8:ex:3}, part
\textbf{f)}). Rebuild, load it, and record what happens to \code{insmod}, and how long it takes
before the console says anything.

\xpart{b} The message that eventually appears mentions neither interrupts nor your driver. Quote it,
and explain the connection.

\xpart{c} This produces no failing test. \code{make test L=L08} reports something other than a
failure; what, and why does the harness draw that distinction?

\xpart{d} From another machine you cannot get, but suppose you could run one command on this one
before it died. Which, and what would it have shown?

\xpart{e} \code{qa-dev} is level-triggered. Would the same missing acknowledgement wedge an
edge-triggered device? What would happen instead?

% ------------------------------------------------------------------------------------------
\exercise[fillaccent]{Interrupts counted two ways}{Cross-check}
\label{c8:ex:9}

Predict a count from the period, measure it, and account for a shortfall that is larger than
measurement error and entirely explicable.

\xpart{a} \textbf{Predict.} Your device is programmed with \code{PERIOD_NS=1000000}. How many
interrupts should it raise in exactly two seconds? Write the number down.

\xpart{b} \textbf{Measure.} Load your module, then take a \textbf{delta} across a measured interval:

\begin{shell}
n0=$(grep qa_irq /proc/interrupts | awk '{ print $2 + $3 }')
t0=$(cut -d' ' -f1 /proc/uptime)
sleep 2
n1=$(grep qa_irq /proc/interrupts | awk '{ print $2 + $3 }')
t1=$(cut -d' ' -f1 /proc/uptime)
\end{shell}

\noindent Compute the actual elapsed time and the actual count. Do this three times.

\xpart{c} \textbf{Reconcile.} The measured count is short by something like 15\%, consistently.
Answer each:
\begin{itemize}
  \item Why a delta rather than reading the count once? What did \secref{c8:app:a2} say happens
    across a reload?
  \item Your elapsed time is not 2.00 s. Does using the measured elapsed time rather than 2.00
    remove the shortfall? By how much does it move it?
  \item The remaining shortfall is not noise: three runs agree to within a few percent of each
    other. So it is systematic. \textbf{Before reading further, write down a hypothesis.}
\end{itemize}

\xpart{d} \textbf{Find out.} The device is not a black box; its source is
\file{tools/qa-dev/qa-dev.c}. Find \code{qa_dev_arm_timer} and read the one line that decides when
the next interrupt happens.
\begin{itemize}
  \item Is the next interrupt scheduled relative to \emph{when the last one fired}, or to an
    absolute schedule?
  \item Given that, what is the actual interval between two interrupts, in terms of the programmed
    period and anything else?
  \item Does that account for a 15\% shortfall at a 1 ms period? Estimate the missing term and say
    whether it is a plausible number for this system.
\end{itemize}

\xpart{e} \textbf{Test the explanation, and watch it fail.} A relative re-arm means each interval is
\code{period + overhead}, so the model predicts a \textbf{constant} overhead: a large shortfall at
short periods and a small one at long ones.

Run it at \code{period_ns=10000000} (10 ms), \code{1000000} (1 ms) and \code{100000} (0.1 ms). For
each, compute the shortfall \emph{and} the overhead your model implies, which is \code{(elapsed /
count) - period}.

The measured result on this target:

\begin{narrowtable}{@{}lrrr@{}}
  \thead{Period} & \thead{Count in \textasciitilde{}2.1 s} & \thead{Shortfall} & \thead{Implied overhead} \\
  \midrule
  10 ms & 190 & 8.7\% & 947 us \\
  1 ms & 1731 & 17.6\% & 213 us \\
  0.1 ms & 13085 & 38.6\% & 63 us \\
\end{narrowtable}

\noindent \textbf{The implied overhead is not constant. It is larger at longer periods, which is
backwards.} The relative re-arm is real and is part of the answer, and it is plainly not the whole
answer.

Do not resolve this by inventing a mechanism. Instead, answer:
\begin{itemize}
  \item What would the table look like if a fixed per-interrupt overhead were the only effect?
  \item Your driver is one of at least three things in the measurement: the driver, the kernel's
    timer and scheduler, and QEMU. Which of the three does your \file{/proc/uptime} measurement
    include?
  \item At 10 ms the guest is idle almost all the time and at 0.1 ms it is not. Why might
    \emph{that} affect how promptly an emulated timer is serviced?
  \item What experiment would separate the device's behaviour from the emulator's? Name one you
    could run, and one you could not run here at all.
\end{itemize}

\xpart{f} \textbf{The fast end is a different effect.} At 0.1 ms compare the driver's own two
counters:

\begin{console}
hard=13708  samples=14389
\end{console}

\noindent More samples arrived than interrupts were taken.
\begin{itemize}
  \item How can one interrupt deliver more than one sample? Answer in terms of level triggering and
    when the line is re-asserted.
  \item At 1 ms the same two counters are equal. What changed?
  \item Which status bit would tell you the device had given up and started dropping samples, and
    did it get set?
\end{itemize}

\xpart{g} A colleague reports ``the device is dropping 15\% of its interrupts, the driver is
broken''. There are at least three things wrong with that sentence. Rewrite it so that every claim
in it is supported by a measurement you have actually taken, and say what you would still not know.

\xpart{h} You are asked for ``the interrupt latency of this device'' for a datasheet. Given
everything above, say what you would put, what conditions you would state alongside it, and what you
would refuse to claim. Chapter~\ref{c12:ch} is the chapter that measures this properly, using a
register that does not exist on real hardware; say why that register is necessary.
